\section{Construction principles and further questions}

The point counts follow from specific changes in finite structure. In the
code constructions, an exchange is governed by the union of conflicts;
allowing signed local replacements enlarges the set of available changes.
In the equatorial construction, moving to a different shell supplies directions
that remain compatible with all existing caps. In the section constructions,
the ambient embedding determines projection multiplicities that are invisible
in the abstract Gram type alone.

These observations suggest concrete mathematical variables for further work.
The support family and sign assignments can change together. Head classes and
tail labels can be chosen by their net contribution to the whole configuration.
An embedded section can be studied through a larger parent root system.
A moment identity can be selected for the statistic that is easiest to witness,
rather than for a full reconstruction of every fibre.

The accompanying finite data preserve these choices as explicit objects.
They make it possible to revisit a construction, test a new local change,
and determine its effect on the final count without repeating the entire
discovery search.
